\documentclass[12pt,a4paper]{ctexart}
\usepackage{geometry}
\geometry{left=2.5cm,right=2.5cm,top=2.5cm,bottom=2.5cm}
\usepackage{listings}
\usepackage{xcolor}
\usepackage{booktabs}
\usepackage{amsmath}
\usepackage{graphicx}

\lstset{
  basicstyle=\small\ttfamily,
  keywordstyle=\color{blue}\bfseries,
  commentstyle=\color{gray}\itshape,
  numbers=left,numberstyle=\tiny\color{gray},
  frame=single,rulecolor=\color{black!30},
  breaklines=true,breakatwhitespace=true,
  language=C++
}

\title{\textbf{编译器设计专题实验报告}\\[0.5em]实验四：SLR(1)分析表生成}
\author{张凌曜 \quad 2234214014 \quad 计算机2306}
\date{}

\begin{document}
\maketitle

\section{实验目的}
在实验三LR(0)规范族的基础上，算FIRST集和FOLLOW集，生成SLR(1)分析表（ACTION和GOTO），利用FOLLOW集解决LR(0)中移进-归约冲突。

\section{实验原理}
LR(0)遇到归约项目就无脑归约，SLR(1)改进了一点：只有当前输入符号在$\text{FOLLOW}(A)$里才归约，这样能消掉不少冲突。

FIRST集就是从某个符号串能推导出来的第一个终结符。FOLLOW集是在某些句型中能紧跟在非终结符后面的终结符。FOLLOW集的计算要迭代，直到集合不再变化。

\section{实验过程}

\subsection{FIRST/FOLLOW集计算}
两个集合都用迭代法，反复扫所有产生式直到不动了为止。FOLLOW集的规则有三条：起始符号的FOLLOW加\$；$A \to \alpha B\beta$时把$\text{FIRST}(\beta)$里非$\epsilon$的元素加到$\text{FOLLOW}(B)$；如果$\beta$能推出$\epsilon$就把$\text{FOLLOW}(A)$也加到$\text{FOLLOW}(B)$。

\subsection{分析表构建}
对每个状态$I_i$里的每个项目：点后面是终结符$a$就填移进$s_j$；归约项目$A \to \alpha\cdot$就在FOLLOW(A)对应的列填归约$r_k$；$S' \to S\cdot$在\$列填acc。

\section{测试结果}

\subsection{四则运算文法}

FOLLOW集结果：

\begin{center}
\begin{tabular}{ll}
$\text{FOLLOW}(E) = \{\$,\ ),\ +\}$ & $\text{FOLLOW}(E') = \{\$\}$ \\
$\text{FOLLOW}(T) = \{\$,\ ),\ *,\ +\}$ & $\text{FOLLOW}(F) = \{\$,\ ),\ *,\ +\}$
\end{tabular}
\end{center}

ACTION表（部分）：

\begin{table}[h]
\centering
\begin{tabular}{lcccccc}
\toprule
状态 & ( & ) & * & + & id & \$ \\
\midrule
$I_0$ & s1 & & & & s5 & \\
$I_2$ & & & & s7 & & acc \\
$I_3$ & & r4 & r4 & r4 & & r4 \\
$I_4$ & & r2 & s8 & r2 & & r2 \\
$I_5$ & & r6 & r6 & r6 & & r6 \\
\bottomrule
\end{tabular}
\end{table}

$I_2$里$E' \to E\cdot$是归约，$E \to E\cdot+T$是移进。LR(0)下这是冲突，但SLR(1)里归约只在$\text{FOLLOW}(E')=\{\$\}$时执行，移进在$+$时执行，两边不重叠，\textbf{冲突解决了}。

\begin{center}
\includegraphics[height=7cm]{image/1.1.png}
\vspace{0.3em}
\includegraphics[height=5cm]{image/1.2.png}
\end{center}

\subsection{嵌套文法}
$S \to aSb | aS | c$

$I_4$里有$S \to aS\cdot b$（移进）和$S \to aS\cdot$（归约）。SLR(1)看$\text{FOLLOW}(S) = \{\$, b\}$，$b$在FOLLOW里所以移进和归约对$b$都有定义，冲突没法完全消掉。程序里对$b$优先移进，并输出了冲突提示。

\begin{center}
\includegraphics[height=7cm]{image/2.png}
\end{center}

\section{选做：flex/bison自动化工具}
用flex做词法分析、bison做语法分析，实现了一个四则运算计算器。

\subsection{flex词法定义（calc.l）}
flex用正则表达式定义Token规则，自动生成了扫描器\texttt{lex.yy.c}：
\begin{lstlisting}[basicstyle=\footnotesize\ttfamily,language={}]
[ \t\r\n]+        { /* skip whitespace */ }
[0-9]+            { yylval = atoi(yytext); return NUM; }
"+"               { return PLUS; }
"*"               { return TIMES; }
"("               { return LPAREN; }
")"               { return RPAREN; }
\end{lstlisting}

\subsection{bison语法定义（calc.y）}
bison用BNF定义文法和语义动作，自动生成了LALR(1)分析器\texttt{calc.tab.c}：
\begin{lstlisting}[basicstyle=\footnotesize\ttfamily,language={}]
%left PLUS MINUS
%left TIMES DIVIDE

expr:
    NUM                { $$ = $1; }
  | expr PLUS expr     { $$ = $1 + $3; }
  | expr TIMES expr    { $$ = $1 * $3; }
  | LPAREN expr RPAREN { $$ = $2; }
  ;
\end{lstlisting}

\texttt{\%left}声明自动解决了运算符优先级和结合性冲突，不需要手动构造分析表。

\subsection{测试结果}
输入\texttt{3 + 5 * 2}，bison按优先级先归约\texttt{5 * 2 = 10}，再归约\texttt{3 + 10 = 13}：

\begin{center}
\includegraphics[height=5cm]{image/3.png}
\end{center}

\textbf{对比手写的SLR(1)程序}：flex/bison把词法分析和语法分析的代码生成全自动化了，只需要写\texttt{.l}和\texttt{.y}两个定义文件，就能得到完整的分析器。这就是工业界用yacc/bison构建编译器的原因。

\section{实验总结}
SLR(1)本质上就是``归约的时候往前看一眼''，用FOLLOW集判断该不该归约。大部分常见冲突都能解决，但像嵌套文法那种就不行了，得上LALR(1)或者LR(1)。

FOLLOW集的迭代计算要注意收敛条件，必须集合完全不变才停，不然会死循环。

选做部分用flex/bison体验了一下自动化工具，发现它们和手写SLR(1)的核心思路是一样的——都是基于LALR(1)分析，只不过flex/bison自动帮你做了FIRST/FOLLOW集计算和冲突解决。这让人更能体会到理论到工程的桥梁作用。

\appendix
\section{附录：代码仓库}
本实验的全部源码、报告及Web演示平台已托管至GitHub：

\texttt{https://github.com/yaossss123/compilation-web}

\end{document}
